% Four-zone theorem dependency skeleton.
% Build dependency-ledger.md and dependency-edges.csv before editing.
% Logical direction: prerequisite -> result proved using it.
\documentclass[UTF8,12pt,tikz,border=9pt]{standalone}
\usepackage{ctex}
\usepackage{amsmath,amssymb}
\usetikzlibrary{arrows.meta,bending,backgrounds,calc,positioning,shapes.geometric}

\definecolor{ink}{HTML}{263548}
\definecolor{muted}{HTML}{6C7888}
\definecolor{foundation}{HTML}{65758B}
\definecolor{routeA}{HTML}{356FA3}
\definecolor{routeAbg}{HTML}{EAF3FA}
\definecolor{routeB}{HTML}{2F7D66}
\definecolor{routeBbg}{HTML}{E8F4EF}
\definecolor{application}{HTML}{A56A16}
\definecolor{applicationbg}{HTML}{FBF2DF}
\definecolor{conclusion}{HTML}{A74343}
\definecolor{conclusionbg}{HTML}{F9EAEA}
\definecolor{alternative}{HTML}{72538F}
\definecolor{contextbg}{HTML}{F2F4F6}

\pgfdeclarelayer{background}
\pgfsetlayers{background,main}

\tikzset{
  >={Stealth[length=5pt 1.1,width'=0pt .55]},
  result/.style={rectangle,rounded corners=2.5pt,draw=ink!72,fill=white,
    line width=.65pt,align=center,inner xsep=5pt,inner ysep=4pt,
    minimum width=3.7cm,minimum height=.9cm,font=\footnotesize,text=ink},
  foundation/.style={result,draw=foundation,fill=contextbg},
  routeA/.style={result,draw=routeA,fill=routeAbg,line width=.9pt},
  routeB/.style={result,draw=routeB,fill=routeBbg,line width=.9pt},
  application/.style={result,draw=application,fill=applicationbg,minimum width=3.15cm},
  conclusion/.style={result,draw=conclusion,fill=conclusionbg,
    minimum width=4.4cm,minimum height=1.25cm,line width=1.15pt,font=\small},
  context/.style={result,draw=foundation!60,minimum height=.72cm,font=\scriptsize},
  gate/.style={diamond,aspect=2.1,draw=application,fill=applicationbg,
    align=center,font=\scriptsize,inner sep=2pt,text=ink},
  direct/.style={->,draw=ink!78,line width=.85pt,shorten >=2pt,shorten <=2pt},
  input/.style={->,draw=muted,line width=.62pt,shorten >=2pt,shorten <=2pt},
  alt/.style={->,draw=alternative,line width=1pt,double distance=.9pt,
    dash pattern=on 4pt off 2.4pt,shorten >=2pt,shorten <=2pt},
  imply/.style={->,draw=muted,line width=.58pt,densely dotted,
    shorten >=2pt,shorten <=2pt},
  zonetitle/.style={font=\small\bfseries,text=ink,anchor=west},
  legend/.style={font=\scriptsize,text=muted,anchor=west},
}

\newcommand{\zone}[6]{%
  \begin{pgfonlayer}{background}
    \fill[#5,rounded corners=5pt] (#1,#2) rectangle (#3,#4);
    \draw[#6,line width=.65pt,rounded corners=5pt] (#1,#2) rectangle (#3,#4);
  \end{pgfonlayer}}

\begin{document}
\begin{tikzpicture}[x=1cm,y=1cm]
  \zone{0}{1.7}{4.8}{13.9}{contextbg}{foundation!60}
  \zone{5.3}{1.7}{14.7}{13.9}{routeAbg!42}{routeA!65}
  \zone{15.2}{1.7}{22.7}{13.9}{applicationbg!48}{application!65}
  \zone{23.2}{1.7}{28.8}{13.9}{conclusionbg!48}{conclusion!65}
  \zone{0}{-.1}{28.8}{1.05}{contextbg}{foundation!60}

  \node[font=\Large\bfseries,text=ink] at (14.4,14.75)
    {主要定理与命题的逻辑依赖图};
  \node[zonetitle] at (.3,13.5) {I. 基础工具};
  \node[zonetitle] at (5.6,13.5) {II. 两条主证明路线};
  \node[zonetitle] at (15.5,13.5) {III. 应用构造与闭合};
  \node[zonetitle] at (23.5,13.5) {IV. 最终结论};

  \node[foundation] (f1) at (2.4,11.7) {基础引理组 A\\引理 1.1、1.2};
  \node[foundation] (f2) at (2.4,9.8) {基础命题 B\\命题 2.3};
  \node[foundation] (f3) at (2.4,6.5) {技术输入 C\\引理 3.4};

  \node[routeA] (a1) at (7.8,11.7) {主定理 A\\定理 4.1};
  \node[routeA] (a2) at (12.2,11.7) {增强/推论 A$'$\\推论 4.5};
  \node[routeB] (b1) at (7.8,7.4) {替代主定理 B\\定理 5.1};
  \node[routeB] (b2) at (12.2,7.4) {优化版本 B$'$\\定理 5.6};

  \node[application] (c1) at (17.25,7.1) {反证假设导出构造\\命题 6.2};
  \node[application] (c2) at (20.85,7.1) {线性无关/非零性\\引理 7.1};
  \node[gate] (g) at (20.0,11.25) {矛盾门\\维数或高度冲突};

  \coordinate (join) at (22.95,10.8);
  \node[conclusion] (t) at (26.0,10.8) {最终定理\\定理 A};
  \node[conclusion] (u) at (26.0,6.3) {第二结论\\定理 C};

  \node[context] (ctx1) at (8.5,.48) {背景比较：经典定理 X};
  \node[context] (ctx2) at (19.7,.48) {补充推广：定理 Y（不参与主证明）};

  \draw[input] (f1.east) -- (a1.west);
  \draw[input,rounded corners=4pt] (f2.east) -- (5.05,9.8) -- (5.05,7.4) -- (b1.west);
  \draw[input,rounded corners=4pt] (f3.east) -- (15.0,6.5) -- (c1.west);
  \draw[imply] (a1) -- (a2);
  \draw[imply] (b1) -- (b2);
  \draw[direct] (c1) -- (c2);

  \draw[alt,rounded corners=4pt] (a2.east) -- (14.95,11.7) -- (14.95,12.7) -- (20.0,12.7) -- (g.north);
  \draw[alt,rounded corners=4pt] (b2.east) -- (15.0,7.4) -- (15.0,9.9) -- (18.9,9.9) -- (g.south west);

  \draw[draw=ink!78,line width=.85pt,rounded corners=4pt] (g.east) -- (22.65,11.25) -- (join);
  \draw[draw=ink!78,line width=.85pt] (c2.east) to[out=0,in=235] (join);
  \draw[direct] (join) -- (t.west);
  \draw[direct,rounded corners=4pt] (join) -- (22.95,6.3) -- (u.west);
  \fill[application] (join) circle (1.6pt);

  \node[zonetitle] at (.3,.62) {补充与比较};
  \draw[direct] (23.6,2.7) -- ++(1.0,0) node[legend] {直接依赖};
  \draw[alt] (23.6,2.25) -- ++(1.0,0) node[legend] {可替代路线};
  \draw[imply] (26.1,2.7) -- ++(1.0,0) node[legend] {推论/推广};
  \draw[input] (26.1,2.25) -- ++(1.0,0) node[legend] {技术输入};
\end{tikzpicture}
\end{document}
